\documentclass[]{../../../tex/classes/styledArticle}

\usepackage{upgreek}
%\usepackage{txfonts}
\usepackage{countriesofeurope}
\input pdfmsym
%\newcommand \bs \blacksquare

\usepackage{titlesec}
\titleformat{\section}[block]
{\huge}
{\dotfill (\textbf{\thesection})\dotfill}
{0em}
{}


\author{\textit{Shahar Perets}}
\title{Module G Practice 2}
\begin{document}
	\maketitle
	
	\large
	
	\section{}
	\textbf{Question:}
	
	Here are your 5 Band III words. Make them feel real — not just defined! Look up the definitions if you don’t know them.
	\begin{center}
		\begin{multicols}{5}
			Inevitable  \\
			Perceive    \\
			Contradict  \\
			Deteriorate \\
			Ambiguous
		\end{multicols}
	\end{center}
	Your tasks:
	\begin{itemize}
		\item Write a sentence for each word that clearly shows its meaning
		\item Which word best describes how you feel this week? Explain why.
		\item Challenge: use 2 of the words in the same sentence
	\end{itemize}
	
	\textbf{Answer: }
	\begin{itemize}
		\item 
		\begin{enumerate}
			\item Once the war started, it was inevitable to cancel the operation.  
			\item I cannot perceive how tall the Burj Khalifa is. 
			\item Emotionalism contradicts any kind of logic. 
			\item Khamenei's regime deteriorated until he was killed. 
			\item Most of the questions I've been asked recently are ambiguous and not well-defined. 
		\end{enumerate}
		\item I think the word ``ambiguous describes my week best, since last week has been full of ambiguous news about the recent events. 
		\item It's imperceivable how many argue, just because of the use of ambiguous definitions. 
	\end{itemize}
	
	\section{}
	\textbf{Question: }
	
	Choose ONE statement below. Do you agree, disagree, or is it complicated?
	\begin{enumerate}
		\item Social media has done more harm than good to our generation.
		\item It is more important to be kind than to be honest.
		\item School prepares you for exams, not for life.
	\end{enumerate}
	Your tasks:
	\begin{itemize}
		\item Write your position in 2–3 sentences
		\item Give one strong reason and one example
		\item Write one counterargument — then respond to it
	\end{itemize}
	\textbf{Answer: }
	\begin{itemize}
		\item \begin{enumerate}
			\item I think social media has done more good than harm, ergo I don't agree with that sentence. Social media can surpass borders and connect large groups of people, which enables new communities to spawn and faster globalization. 
			\item I think that it's more important to be honest than kind, hence I don't agree with the second sentence either. There is no point in ``fake'' kindness -- I'll most likely result, in either an ``expulsion'' of interests in the future, or cause one to believe more wrong behaviors, as they're reinforced by dishonest kindness from his environment. 
			\item School does prepare you for exams, and not for life. Consider even the content of these exams -- a math exam is far away from what is being learned in universities, those who finished computer science exams don't know how to code in practice, some finish English finals and don't know how to talk in English, etc. Add to that the fact the ``real'' life involves much more than just pure knowledge; success requires abilities to lead, to manage, to work within a group and more, which in my experience, school fails to give. 
		\end{enumerate}
		\item Already when I explained my argument. 
		\item \begin{enumerate}
			\item counterargument: Social media is addictive. 
			
			Response: Yes, indeed. But consider a place where, when two people want to talk, they're forced to meet each other? When in a war or a pandemic, you're only going to hear from a limited number of people? Where every large community disappears, and the most efficient way to spread knowledge -- the internet -- is gone? It would be a worse world than the one we live in, therefore, social media does more good than harm, and the existence of harmful effects of social media is not sufficient to prove otherwise. 
			
			\item counterargument: but kindness causes people to be happy. 
			
			Repose: Suppose one person is kind to the other, even though the other greatly disturbs him. He chose not to be honest about the distraction. Yes, the other may feel ``down'' for a bit, if the first had noted the distraction, but overall, apart from a minor inconvenience of being honest yet somewhat unkind, it'll result in a better position where no one bothers the other. 
			
			Hence, being honest may cause unhappiness for a little bit, but in the long run, honesty will result in everyone being happier. 
			
			\item counterargument: Schools have plenty of activities that are not just learning.
			
			Repose: those activities are usually short, not long-term, and in a small group that does not represent big organizations or corporations. 
		\end{enumerate}
	\end{itemize}
	
\end{document}